\PassOptionsToPackage{usenames,dvipsnames}{xcolor}
\documentclass[a4paper,12pt]{letter}
\usepackage[utf8]{inputenc}
\usepackage[T1]{fontenc}
\usepackage[english,slovene]{babel}
\usepackage[tikz]{bclogo}
\usepackage{qrcode,marvosym}
\usepackage{pgf,tikz,pgfplots}
\pgfplotsset{compat=1.14}
\usepackage{mathrsfs}
\usetikzlibrary{arrows}
\usepackage{graphicx}
\newcommand{\ds}{\displaystyle}
\usepackage{fancybox}
\usepackage{amsfonts}
\usepackage{amsmath}
\usepackage{wallpaper}
\usepackage{hyperref}
\usepackage{array}
\newcolumntype{M}[1]{>{\centering\arraybackslash}m{#1}}
\newcolumntype{N}{@{}m{0pt}@{}}
\usepackage{fancyhdr,enumerate}
\newcommand{\ora}{\overrightarrow}
\renewcommand{\headrulewidth}{3pt}
\renewcommand{\headrule}{\hbox to\headwidth{%
  \color{gray}\leaders\hrule height \headrulewidth\hfill}}
\renewcommand{\footrulewidth}{2pt}
\renewcommand{\footrule}{\hbox to\headwidth{%
  \color{brown}\leaders\hrule height \footrulewidth\hfill}}
\usepackage{gensymb}
\usepackage{amssymb}
  \textwidth 20 true cm
  \oddsidemargin -2 true cm
  \evensidemargin -2 true cm
  \topmargin -3.5 true cm
  \textheight 27.5 true cm
  \linespread{1.2}
\renewcommand{\headrulewidth}{0.3pt}
\renewcommand{\footrulewidth}{1.9pt}
\definecolor{qqwuqq}{rgb}{1,0,0}
\definecolor{qqqqff}{rgb}{0,0,1}
\definecolor{qqffqq}{rgb}{0,1,0}
\definecolor{ffqqqq}{rgb}{1,0,0}
\definecolor{zzuxux}{rgb}{0,0,0}
\usepackage{mdframed}
\usepackage{multicol}
\definecolor{bgblue}{RGB}{0,0,253}
\usepackage{eurosym}
%%%%%%%%%%%%%%%%%%%%%%%%%%%%%%%%%%%%%%%%%%%%%%%%%%%%%%%%
\newcounter{tocke}
\setcounter{tocke}{0}
\usepackage{tikz-3dplot}
%%%%%%%%%%%%%%%%%%%%%%%%%%%%%%%%%%%%%%%%%%%%%%%%%%%%%%%%%%
\mdfdefinestyle{theoremstyle}{%
linecolor=brown!40,linewidth=1pt,%
frametitlerule=true,%
frametitlebackgroundcolor=brown!50,
innertopmargin=\topskip,
}
\mdfdefinestyle{theoremstyle1}{%
linecolor=brown,middlelinewidth=1pt,%
frametitlerule=true,%
apptotikzsetting={\tikzset{mdfframetitlebackground/.append style={%
shade,left color=black!40, right color=gray!10},
mdfbackground/.append style={
top color=white!100!white,
bottom color=black!5
},
}},
frametitlerulecolor=gray,
frametitlerulewidth=2pt,
innertopmargin=\topskip,
innerbottommargin=\topskip,
}
\mdtheorem[style=theoremstyle1]{naloga}{Naloga}
\mdtheorem[style=theoremstyle]{resitev}{Analiza Naloge}
\mdtheorem[style=theoremstyle]{kriterij}{\v{S}tevilo dose\v{z}enih to\v{c}k na testu}
%%%%%%%%%%%%%%%%%%%%%%%%%%%%%%%%%%%%%%%%%%%%%%%%%%%%%%%%%%
\usepackage[table]{xcolor}
\usepackage[printwatermark]{xwatermark}
\newwatermark[allpages,color=brown!2,angle=45,scale=5,
xpos=-5,ypos=0]{matej.info}
%%%%%%%%%%%%%%%%%%%%%%%%%%%%%%%%%%%%%%%%%%%%%%%%%%%%%%%%%%

\begin{document}
\pagestyle{fancy}
\renewcommand\headrulewidth{0pt}
\lhead{}\chead{}\rhead{}
\rfoot{\vspace*{-2\baselineskip}}
\renewcommand{\vec}{\overrightarrow}
%%%%%%%%%%%%%%%%%%%% podatki testa -- tu spreminjaj
\def\programletnik{Gim - 3}
\def\naslov{Test 1.2 - }
\def\sklop{Eksponentna in logaritemska funkcija}
%%%%%%%%%%%%%%%%%%%% tocke po nalogah
\def\tockeA{4+3}
\def\tockeB{2+3+4+4}
\def\tockeC{4+4}
\def\tockeD{4+4}
\def\tockeE{4}
%%%%%%%%%%%%%%%%%%%%
\begin{minipage}{20cm}
\begin{bclogo}[couleur=brown!30,
  arrondi=0,ombre=true, couleurOmbre=gray!50,
logo =\bcplume,
  noborder=true]
{\textrm{{\Large{\naslov \sklop} \Huge{}
\large{} \hfill $\programletnik$}}}
\end{bclogo}
\end{minipage}
\begin{minipage}{20cm}
\begin{bclogo}[couleur=brown!30,
  arrondi=0,ombre=true, couleurOmbre=gray!50,
logo =\bcplume,
  noborder=true]
{\textrm{{\raisebox{-0.3cm}{\large{Ime in Priimek: \rule{10cm}{0.1mm}} \Huge{}
\large{} \hfill $ $}}}}
\end{bclogo}
\end{minipage}
%%%%%%%%%%%%%%%%%%%%%%%%%%%%%%%%%%%%%%%%%%%%%%%%%

%%%%%%%%%%%%%%%%%%%%%%%%%%%%%%%%%%%%%%%%%%%%%%%%%
\begin{naloga}[\hfill $\tockeA$] \addtocounter{tocke}{\numexpr\tockeA\relax}
Re\v{s}i ena\v{c}bo:
\begin{enumerate}[a)]
\item $5^{x-1}=5^{x}-25\cdot2^{x-1}$
\item $4^{x}=5$\quad (rezultat zaokro\v{z}i na \v{s}tiri decimalna mesta).
\end{enumerate}
\end{naloga}\newpage

%%%%%%%%%%%%%%%%%%%%%%%%%%%%%%%%%%%%%%%%%%%%%%%%%
\begin{naloga}[\hfill $\tockeB$] \addtocounter{tocke}{\numexpr\tockeB\relax}
Dana je funkcija $f(x)=\log_{1/2}(x+2)-1$.
\vspace{-0.4cm}

\begin{multicols}{2}
\begin{enumerate}[a)]
\item Izra\v{c}unaj $f(14)$.
\item Za kateri $x$ je vrednost funkcije $f$ enaka $-6$?
\item Zapi\v{s}i inverzno funkcijo $f^{-1}$ in jo nari\v{s}i.
\item Izra\v{c}unaj prese\v{c}i\v{s}\v{c}e funkcije $f$ s funkcijo \\$g(x)=\log_2(x-5)-4$.
\end{enumerate}
\end{multicols}
\end{naloga}
\vspace{0.3cm}
\begin{tikzpicture}[scale=0.82]
\draw[very thin,gray!30] (-3,-7) grid (10,7);
\draw[->,thick] (-3,0) -- (10,0) node[right] {$x$};
\draw[->,thick] (0,-7) -- (0,7) node[above] {$y$};
\end{tikzpicture}
\newpage

%%%%%%%%%%%%%%%%%%%%%%%%%%%%%%%%%%%%%%%%%%%%%%%%%
\begin{naloga}[\hfill $\tockeC$] \addtocounter{tocke}{\numexpr\tockeC\relax}
\begin{enumerate}[a)]
\item Izra\v{c}unaj
\[
\log_3 3-\ln e+\log_2\frac{1}{16}+3^{\log_3 2}.
\]
\item Izra\v{c}unaj
\[
\log\frac12+\log\frac23+\log\frac34+\cdots+\log\frac{99}{100}.
\]
\end{enumerate}
\end{naloga}\newpage

%%%%%%%%%%%%%%%%%%%%%%%%%%%%%%%%%%%%%%%%%%%%%%%%%
\begin{naloga}[\hfill $\tockeD$] \addtocounter{tocke}{\numexpr\tockeD\relax}
\begin{enumerate}[a)]
\item Dolo\v{c}i $x$, \v{c}e je
\[
\log x=2+\log_4 20-\log_4 5-\log_4 32.
\]
\item Naj bo $p=\ln n$, $q=\ln m$. Izrazi s $p$ in $q$ vrednost logaritma
\[
\ln\!\left(\frac{m\,n^2}{\sqrt e}\right).
\]
\end{enumerate}
\end{naloga}\newpage

%%%%%%%%%%%%%%%%%%%%%%%%%%%%%%%%%%%%%%%%%%%%%%%%%
\begin{naloga}[\hfill $\tockeE$] \addtocounter{tocke}{\numexpr\tockeE\relax}
Dolo\v{c}i $x$, da bo funkcija
\[
y=\ln\!\left(x^2+x\right)-2\ln\!\left(\sqrt2\right)
\]
pozitivna.
\end{naloga}

\vfill
\begin{kriterij*} \begin{center} \renewcommand{\arraystretch}{2.5}
\begin{tabular}{|>{\columncolor{gray!20}}c||c|c|c|c|c||c|c|c|} \hline Odstotki &$[0,45)$ & $[45,60)$ & $[60,75)$ & $[75,90)$ & $[90,100]$& \v{S}t. mo\v{z}nih to\v{c}k& Dose\v{z}ene to\v{c}ke & OCENA \\ \hline Ocena & 1 & 2 & 3 & 4 & 5 &\cellcolor{gray!20} $\thetocke$&\cellcolor{gray!20} \phantom{ocena} & \cellcolor{gray!20} \phantom{a} \\ \hline \end{tabular} \end{center} \vspace{0.3cm}  \end{kriterij*}

\newpage
\textsc{\Large Re\v{s}itve s postopki}
\vspace{0.5cm}

\begin{resitev}[-- eksponentni ena\v{c}bi]
\textbf{a)} $5^{x-1}=5^{x}-25\cdot2^{x-1}$. Ker je $5^{x-1}=\dfrac{5^x}{5}$, prenesemo vse \v{c}lene z $5^x$ na eno stran:
\[
25\cdot2^{x-1}=5^x-\frac{5^x}{5}=\frac{4}{5}\cdot5^x.
\]
Pomno\v{z}imo z $5$ in zapi\v{s}emo $2^{x-1}=\dfrac{2^x}{2}$:
\[
125\cdot\frac{2^x}{2}=4\cdot5^x\ \Longrightarrow\ \frac{125}{2}\cdot2^x=4\cdot5^x.
\]
Delimo z $5^x$ in z $\dfrac{125}{2}$:
\[
\left(\frac25\right)^x=\frac{4\cdot2}{125}=\frac{8}{125}=\left(\frac25\right)^3\ \Longrightarrow\ x=3.
\]
\textit{Preverimo:} $5^{2}=25$ in $5^3-25\cdot2^2=125-100=25$. \checkmark

\textbf{b)} $4^x=5$. Logaritmiramo:
\[
x=\log_4 5=\frac{\ln5}{\ln4}=\frac{\log5}{\log4}\approx\frac{0{,}6990}{0{,}6021}\approx1{,}1610.
\]
\end{resitev}

\begin{resitev}[-- logaritemska funkcija in njen inverz]
Dana je $f(x)=\log_{1/2}(x+2)-1$, definicijsko obmo\v{c}je $x>-2$ (navpi\v{c}na asimptota $x=-2$). Ker je osnova $\tfrac12<1$, je funkcija \textbf{padajo\v{c}a}.

\textbf{a)}
\[
f(14)=\log_{1/2}(16)-1=\log_{1/2}\!\left(2^4\right)-1=-4-1=-5.
\]

\textbf{b)} Re\v{s}imo $f(x)=-6$:
\[
\log_{1/2}(x+2)-1=-6\ \Longrightarrow\ \log_{1/2}(x+2)=-5\ \Longrightarrow\ x+2=\left(\frac12\right)^{-5}=32\ \Longrightarrow\ x=30.
\]

\textbf{c)} Zamenjamo vlogi $x$ in $y$ ter izrazimo nov $y$:
\[
y=\log_{1/2}(x+2)-1\ \Longrightarrow\ x=\log_{1/2}(y+2)-1\ \Longrightarrow\ x+1=\log_{1/2}(y+2)
\]
\[
\Longrightarrow\ y+2=\left(\frac12\right)^{x+1}\ \Longrightarrow\ y=\left(\frac12\right)^{x+1}-2.
\]
Torej
\[
f^{-1}(x)=\left(\frac12\right)^{x+1}-2,\qquad D_{f^{-1}}=\mathbb R,\qquad \text{horizontalna asimptota } y=-2.
\]
Graf $f^{-1}$ je zrcalna slika grafa $f$ prek premice $y=x$ (obe krivulji se sekata na tej premici v to\v{c}ki $(-1,-1)$, saj je $f(-1)=-1$).

\begin{center}
\begin{tikzpicture}[scale=0.5]
\draw[very thin,gray!25] (-5,-6) grid (11,10);
\draw[->,thick] (-5,0) -- (11,0) node[right] {$x$};
\draw[->,thick] (0,-6) -- (0,10) node[above] {$y$};
\draw[dashed,thick,red!60] (-2,-6) -- (-2,10) node[pos=0.96,right]{\scriptsize $x=-2$};
\draw[dashed,thick,blue!60] (-5,-2) -- (11,-2) node[pos=0.98,above]{\scriptsize $y=-2$};
\draw[dotted] (-5,-5) -- (10,10);
\draw[domain=-1.96:10.9,smooth,variable=\x,red,thick,samples=150] plot ({\x},{-ln(\x+2)/ln(2)-1});
\draw[domain=-4.5:10.9,smooth,variable=\x,blue,thick,samples=150] plot ({\x},{0.5^(\x+1)-2});
\filldraw[black] (-1,-1) circle (2pt) node[above right]{\scriptsize $(-1,-1)$};
\filldraw[red] (6,-4) circle (2pt) node[below right]{\scriptsize $(6,-4)$};
\node[red] at (8,1.2) {$f$};
\node[blue] at (1.2,8) {$f^{-1}$};
\end{tikzpicture}
\end{center}
\emph{(To\v{c}ki $(14,-5)$ in $(30,-6)$ iz a) in b) le\v{z}ita na krivulji $f$ precej dlje v desno/navzdol in zaradi preglednosti nista prikazani na tej skici.)}

\textbf{d)} I\v{s}\v{c}emo presek s $g(x)=\log_2(x-5)-4$ (definicijsko obmo\v{c}je $x>5$). Ker je $\log_{1/2}t=-\log_2 t$, lahko $f$ zapi\v{s}emo kot $f(x)=-\log_2(x+2)-1$. Iz $f(x)=g(x)$ sledi
\[
-\log_2(x+2)-1=\log_2(x-5)-4\ \Longrightarrow\ \log_2(x+2)+\log_2(x-5)=3
\]
\[
\Longrightarrow\ \log_2\big[(x+2)(x-5)\big]=3\ \Longrightarrow\ (x+2)(x-5)=2^3=8.
\]
Uredimo:
\[
x^2-3x-10=8\ \Longrightarrow\ x^2-3x-18=0\ \Longrightarrow\ (x-6)(x+3)=0\ \Longrightarrow\ x=6\ \text{ali}\ x=-3.
\]
Ker mora biti $x>5$ (zaradi $g$), je edina veljavna re\v{s}itev $x=6$. Preverimo:
\[
f(6)=\log_{1/2}(8)-1=-3-1=-4,\qquad g(6)=\log_2(1)-4=0-4=-4.\ \checkmark
\]
Prese\v{c}i\v{s}\v{c}e je torej to\v{c}ka $(6,-4)$.
\end{resitev}

\newpage
\begin{resitev}[-- ra\v{c}unanje z logaritmi]
\textbf{a)} Vsak \v{c}len poenostavimo posebej: $\log_3 3=1$, $\ln e=1$, $\log_2\dfrac{1}{16}=\log_2\!\left(2^{-4}\right)=-4$, in po pravilu $a^{\log_a x}=x$ je $3^{\log_3 2}=2$. Torej
\[
\log_3 3-\ln e+\log_2\frac{1}{16}+3^{\log_3 2}=1-1-4+2=-2.
\]

\textbf{b)} Vsoto logaritmov zapi\v{s}emo kot logaritem produkta:
\[
\log\frac12+\log\frac23+\cdots+\log\frac{99}{100}=\log\!\left(\frac12\cdot\frac23\cdot\frac34\cdots\frac{99}{100}\right).
\]
Produkt je teleskopski -- vsak \v{s}tevec (razen prvega) se okraj\v{s}a s prej\v{s}njim imenovalcem:
\[
\frac12\cdot\frac23\cdot\frac34\cdots\frac{99}{100}=\frac{1}{100}.
\]
Torej je vsota
\[
\log\frac{1}{100}=\log\!\left(10^{-2}\right)=-2.
\]
\end{resitev}

\begin{resitev}[-- ra\v{c}unanje z logaritmi II]
\textbf{a)} Vse logaritme z desne strani (razen konstante $2$) najprej zdru\v{z}imo -- gre za logaritme z osnovo $4$:
\[
\log_4 20-\log_4 5-\log_4 32=\log_4\!\left(\frac{20}{5}\right)-\log_4 32=\log_4 4-\log_4 32=1-\log_4 32.
\]
Ker je $4^{5/2}=4^2\cdot4^{1/2}=16\cdot2=32$, je $\log_4 32=\dfrac52$, zato
\[
\log_4 20-\log_4 5-\log_4 32=1-\frac52=-\frac32.
\]
Torej
\[
\log x=2+\left(-\frac32\right)=\frac12.
\]
Sledi
\[
x=10^{1/2}=\sqrt{10}\approx3{,}162.
\]

\textbf{b)} Izraz razstavimo s pravili za logaritem produkta, kvocienta, potence in korena:
\[
\ln\!\left(\frac{m\,n^2}{\sqrt e}\right)=\ln m+\ln\!\left(n^2\right)-\ln\sqrt e=\ln m+2\ln n-\frac12\ln e.
\]
Ker je $\ln n=p$, $\ln m=q$ in $\ln e=1$, sledi
\[
\ln\!\left(\frac{m\,n^2}{\sqrt e}\right)=q+2p-\frac12.
\]
\end{resitev}
\newpage

\begin{resitev}[-- logaritemska neena\v{c}ba]
Najprej dolo\v{c}imo definicijsko obmo\v{c}je: argument logaritma mora biti pozitiven,
\[
x^2+x>0\ \Longrightarrow\ x(x+1)>0\ \Longrightarrow\ x<-1\ \text{ali}\ x>0.
\]
\v{Z}elimo, da je $y>0$:
\[
\ln\!\left(x^2+x\right)-2\ln\sqrt2>0.
\]
Ker je $2\ln\sqrt2=\ln\!\left(\sqrt2\right)^2=\ln2$, se pogoj poenostavi v
\[
\ln\!\left(x^2+x\right)>\ln2.
\]
Ker je $\ln$ nara\v{s}\v{c}ajo\v{c}a funkcija, je to ekvivalentno
\[
x^2+x>2\ \Longrightarrow\ x^2+x-2>0\ \Longrightarrow\ (x+2)(x-1)>0\ \Longrightarrow\ x<-2\ \text{ali}\ x>1.
\]
To obmo\v{c}je je \v{z}e v celoti vsebovano v definicijskem obmo\v{c}ju ($x<-2\subset x<-1$ in $x>1\subset x>0$), zato je kon\v{c}na re\v{s}itev
\[
x\in(-\infty,-2)\cup(1,\infty).
\]
\end{resitev}

\end{document}

